\section{研究背景}

\subsection{临床与机器人学动机}

手术技术水平直接影响临床结局。一项针对减重外科医生的研究由同行对手术录像评分，评分最低四分位组的并发症、再次手术和再入院率明显高于评分最高四分位组\citep{birkmeyer2013skill}。目前技能评估主要依靠专家观察，并借助 OSATS 等结构化量表和整体评分\citep{martin1997osats}。这一过程耗时、昂贵且带有主观性，因此很难为受训者提供频繁反馈，也难以在资格认证中保持一致\citep{reiley2011review}。机器人辅助微创手术改变了这一局面，因为机器人本身就在记录外科医生的动作：da~Vinci 系统以固定频率记录每个机械臂的位姿、速度和夹钳角度，内窥镜则记录术野画面。客观技能评估因此可以基于本来就会采集的数据完成，这使它成为医疗机器人智能中的一个自然问题。这类模型只有在面对从未见过的外科医生时依然有效，才具有实际价值。被系统评估的受训者不会出现在训练集中，所以评估必须整体留出外科医生，而不是留出单个试次。

\subsection{相关工作}

JHU-ISI 手势与技能评估数据集（JIGSAWS）是该问题的标准公开基准\citep{gao2014jigsaws}。它包含八位外科医生在台式模型上完成缝合、打结和穿针任务时同步记录的运动学数据和立体视频，每个试次附有自报技能等级和改良 OSATS 评分。该数据集的基准协议区分留一超试次（LOSO，每位外科医生各留出一次重复）和留一用户（LOUO，留出某位外科医生的全部试次）两种划分\citep{ahmidi2017benchmarks}。基于运动学的方法大致分为三类。第一类用全局运动统计量概括一个试次，例如完成时间、路径长度、运动经济性和夹钳动作次数，这些指标与专家评分时关注的要素接近，计算代价也低\citep{reiley2011review, lefor2020motion, anh2020features}。第二类刻画运动信号的纹理。本课题的背景文献 \citet{zia2018rmis} 在每个运动学通道上计算近似熵\citep{pincus1991apen}等度量，用最近邻分类器在 LOUO 下报告了缝合任务 88.2\% 的准确率；同一团队也将这一思路用于视频和加速度计数据\citep{zia2018video}。第三类端到端学习特征，例如在原始运动学序列\citep{fawaz2018cnn, wang2018cnn}或视频片段\citep{funke2019video, liu2021unified}上训练卷积网络；近期的视频 Transformer 已能从临床录像中解码外科医生的动作和技能\citep{kiyasseh2023vit, lavanchy2021threestage}。许多最高的已发表结果采用 LOSO，此时每位测试外科医生也出现在训练集中，这种划分下接近 100\% 的准确率很常见\citep{fawaz2018cnn, funke2019video}。在 LOUO 下，任务变成仅凭七位训练外科医生去泛化到新的人，任何估计的不确定性都很大\citep{varoquaux2018cv}。因此两种划分得到的结果不可直接比较。类别分布带来了第二个难点。新手试次占缝合任务的将近一半，准确率因此会偏向倾向多数类的分类器；不平衡学习的文献建议采用类别平衡的训练方式，并使用对每个类别同等加权的指标，例如 macro-F1\citep{he2009imbalanced}。本文同时采用了这两项建议。

\subsection{方法选择的动机}

本课题规定了核心实验：用简单的试次级统计量构建仅运动学基线，加入来自采样帧或短片段的视频特征，评估一个简单的运动学与视频融合模型，并在 LOUO 下以 macro-F1 与多数类基线比较。我们的设计选择源于数据的两个性质。第一，JIGSAWS 缝合任务只有来自八位外科医生的 39 个可用试次，并且技能标签在同一外科医生内部恒定。因此学习“技能等级之间差异”的有效样本量更接近 8 而不是 39，中级和专家两类各只有两位外科医生。这一点更适合参数少、正则化强的低容量模型，而不是端到端深度网络；对视频我们使用冻结的预训练图像特征加线性分类器，不做任何微调。第二，由于标签属于外科医生，任何能识别外科医生身份的特征都能在训练集内预测标签，却并不反映技能。因此我们把捷径检测作为方法的一部分，在必做实验之外加入两项测量：一是各特征块在多大程度上编码了外科医生身份的探针，二是衡量决策中有多少可以由试次节奏解释的对照实验。作为针对性改进，我们检验背景文献的核心观点，即专业水平体现在动作的规律性上，而不仅仅体现在速度上，具体做法是在简单统计量之外加入运动纹理特征。所有模型、指标和比较都在计算第一个结果之前写入书面协议，因此最终报告哪个模型不会受到结果的影响。
